\documentclass[blue]{beamer}
\usetheme{metropolis}
\usepackage{amssymb,latexsym}
\usepackage{amsmath}
\usepackage{amsthm}
\usepackage{graphicx}
\usepackage{subfigure}
\usepackage{hyperref}

\definecolor{Red}{rgb}{1,0,0}
\definecolor{Blue}{rgb}{0,0,1}
\definecolor{darkblue}{rgb}{0,0,.75}
\definecolor{Green}{rgb}{0,1,0}
\definecolor{magenta}{rgb}{1,0,.6}
\definecolor{lightblue}{rgb}{0,.5,1}
\definecolor{lightpurple}{rgb}{.6,.4,1}
\definecolor{gold}{rgb}{.6,.5,0}
\definecolor{orange}{rgb}{1,0.4,0}
\definecolor{hotpink}{rgb}{1,0,0.5}
\definecolor{newcolor2}{rgb}{.5,.3,.5}
\definecolor{newcolor}{rgb}{0,.3,1}
\definecolor{newcolor3}{rgb}{1,0,.35}
\definecolor{darkgreen1}{rgb}{0, .35, 0}
\definecolor{darkgreen}{rgb}{0, .6, 0}
\definecolor{darkred}{rgb}{.75,0,0}
%\usepackage{beamerthemesplit}
\usepackage{color}
\usepackage{multirow}

\usepackage{lipsum}
\usefonttheme{professionalfonts}
\newcommand\blfootnote[1]{%
  \begingroup
  \renewcommand\thefootnote{}\footnote{#1}%
  \addtocounter{footnote}{-1}%
  \endgroup
}
\newcommand{\E}{\mathbb{E}}
\newcommand{\bi}{\begin{itemize}}
\newcommand{\ei}{\end{itemize}}
%\AtBeginSection{}
\setbeamertemplate{section in toc}[sections numbered]
\setbeamerfont{itemize/enumerate body}{size=\normalsize}
\setbeamerfont{itemize/enumerate subbody}{size=\normalsize}

%\usepackage{amsmath,amssymb,amsthm}
%\usepackage{graphicx}
\usepackage{booktabs}
\usepackage{tikz}
\usetikzlibrary{arrows.meta,positioning,decorations.pathreplacing}
%\usepackage{hyperref}


\title[Chapter1]{Chapter 1: The Subject Matter}
\author[]{}

\date[May 2026]{May 2026}





\begin{document}
  \frame
  {
    \titlepage
  }


% ============================================================
% OUTLINE
% ============================================================
\begin{frame}{Outline}
	\tableofcontents
\end{frame}

% ============================================================
\section{What Is Macroeconomics?}
% ============================================================

\begin{frame}{Defining the Field}
	Several definitions of macroeconomics, each with caveats:
	\vspace{-5pt}
	\begin{enumerate}
		\item \textbf{Study of aggregates}: GDP, inflation, unemployment, interest rates
		\begin{itemize}
			\item But: macroeconomics today is deeply concerned with \emph{distributions} (income, wealth), not just aggregates
		\end{itemize}
		\item \textbf{Quantitative study of general equilibrium}
		\begin{itemize}
			\item But: full GE is not always necessary; single-country analysis is inherently partial
			\item Key insight: behavior of consumers and firms given prices often matters more than market clearing
		\end{itemize}
		\item \textbf{Study of dynamics}: investment, expectations, transitions
		\begin{itemize}
			\item But: static analysis can shed fundamental light too
		\end{itemize}
	\end{enumerate}
		\vspace{-5pt}
	\textbf{Additionally}: Macroeconomics tracks current events---the field follows the major issues of the time (Great Recession, COVID-19, 2025 tariffs).
\end{frame}

\begin{frame}{Goals of This Chapter}
	Two main goals:
	
	\bigskip
	\begin{enumerate}
		\item Introduce the \textbf{central questions} that macroeconomists study
		\item Emphasize a key theme: macroeconomics is fundamentally an \textbf{empirical} field
		\begin{itemize}
			\item \textbf{Measurement} and \textbf{theory} are intertwined
			\item They have evolved in tandem: new data inspires new theory, and new theory creates demand for new data
		\end{itemize}
	\end{enumerate}
	
	\bigskip
	\textbf{Approach}: Walk through important events in the development of the field, illustrating how measurement and theory responded to the major needs of the time.
	
	\bigskip
	The chapter also motivates the empirical and theoretical {methods} used throughout the textbook.
\end{frame}


% ============================================================
\section{The Great Depression: What Is Going On?}
% ============================================================

\begin{frame}{The Birth of National Accounts}
	About 100 years ago: macroeconomists had virtually \textbf{no systematic data}---only scattered accounts of production figures and prices.
	
	\textbf{The Great Depression} made systematic measurement urgent:
	\begin{itemize}
		\item National income and product accounts began in the 1920s
		\item Major research achievements: Colin Clark, Simon Kuznets, Richard Stone
		\item Newly available data helped Keynes and others analyze the macroeconomy
	\end{itemize}
	\textbf{Keynes}: Emphasized market imperfections (sluggish wages and prices), leaving room for demand declines to affect production. Simple theories connecting output, consumption, and investment were developed and remain influential today.
\end{frame}


\begin{frame}{What Do We Measure? GDP and Price Indices}
	\textbf{GDP}: Measure of how much the market economy is producing.
		\vspace{-6pt}
	\begin{itemize}
		\item Does not count leisure, amenities, illegal transactions
		\item Can be measured three ways: sum of final expenditures, sum of incomes, sum of value added
		\item Income splits into labor income (wages/salaries) and capital income (profits, interest, dividends)
	\end{itemize}
	\vspace{-6pt}
	\textbf{Price indices}: Needed to separate nominal changes into quantities and prices.
		\vspace{-5pt}
	\begin{itemize}
		\item Weighted averages of prices. Weights = expenditure shares
		\item Choices: arithmetic vs.\ geometric; base-year vs.\ evolving weights (e.g., T\"ornqvist)
		\item Often guided by microeconomic theory 
	\end{itemize}
	
	\vspace{-5pt}
	\textbf{Capital treatment}: Expenditures on capital are final expenditures (counted in GDP), though capital is in some sense an intermediate good used in future production.
\end{frame}

\begin{frame}{The Pandemic: ``What Is Going On?'' Redux}
	A parallel to the Great Depression occurred during {COVID-19}:
	\begin{itemize}
		\item Economies changed massively overnight
		\item Economists needed \textbf{high-frequency, real-time} measures of activity
		\item Used transaction and commercial data (e.g., credit card spending patterns)
		\item Credit card data also valuable for measuring marginal propensities to consume out of government transfers
	\end{itemize}
	
	\bigskip
	\textbf{Lesson}: ``What is going on?'' is a question that recurs whenever the economy is hit by unprecedented shocks. New data needs arise each time.
\end{frame}

% ============================================================
\section{Long-Run Growth}
% ============================================================

\begin{frame}{Solow and Growth Accounting}
	By the mid-1950s: enough data to analyze growth performance.
	

	\textbf{Solow's two papers}:
		\vspace{-6pt}
	\begin{itemize}
		\item 1956: The neoclassical growth model
		\item 1957: Growth accounting---decompose aggregate growth into contributions of each input plus a residual (``Solow residual'' = technical change)
	\end{itemize}
	

	\textbf{Key finding}: Technical change is of great, and direct, importance for US growth.
	

	\textbf{Reaction}: Was the residual overstated?
		\vspace{-6pt}
	\begin{itemize}
		\item Workers' skills improving (schooling, on-the-job learning)
		\item Mincer (1958): one more year of schooling $\approx$ 10\% higher wage $\Rightarrow$ ``human capital''
		\item Hall and Jorgenson (1969): user cost of capital---measuring the year-by-year price of capital
	\end{itemize}
\end{frame}

\begin{frame}{Cross-Country Comparisons and Endogenous Growth}
	\textbf{Measurement across countries}:
	\begin{itemize}
		\item Kravis, Summers, and Heston (1978): comparable data for 100 countries $\Rightarrow$ Penn World Tables
		\item Maddison: GDP estimates going back to early 19th century
		\item PPP adjustment: purchasing power varies across countries
	\end{itemize}
	
	\bigskip
	\textbf{Key question}: What makes some countries grow fast and others stagnate?
	
	\bigskip
	\textbf{Endogenous growth literature} (late 1980s--1990s):
	\begin{itemize}
		\item Endogenous nature of technological change (incentives to innovate)
		\item Human capital accumulation
		\item Patent data: Who becomes an inventor?
	\end{itemize}
\end{frame}


\begin{frame}{Structural Change}
	\textbf{Structural change}: Some sectors shrink, others grow.
		\vspace{-6pt}
	\begin{itemize}
		\item Typical development: agriculture $\to$ manufacturing $\to$ services
		\item Agriculture: $\sim$1\% of US workforce vs.\ 80\% in poorest countries
		\item Manufacturing overtaken by services
	\end{itemize}

	\textbf{Current concerns}:
		\vspace{-6pt}
	\begin{itemize}
		\item Automation and AI: implications for performance, inequality, firm competition
		\item Women's labor-force participation: steady rise in US postwar 
		\item Government sector size and role
		\item International trade patterns
	\end{itemize}
	

	These are slow-moving trends, never causing immediate media attention, but crucial for economic welfare. Research on growth has come in waves.
\end{frame}

% ============================================================
\section{The 1970s: Stagflation and the Lucas Critique}
% ============================================================

\begin{frame}{The 1970s: An ``Oops''}
	Post-war period to early 1970s: steady growth, Keynesian ``fine tuning.''
	
	\bigskip
	\textbf{Then}: The 1973 oil crisis (OPEC) and a sharp recession.
	\begin{itemize}
		\item \textbf{Productivity slowdown}: appeared permanent to many
		\item \textbf{Rising inequality}: different expressions in different countries
		\begin{itemize}
			\item US and UK: sharply rising wage inequality
			\item EU countries: very high and persistent unemployment
		\end{itemize}
		\item Expansionary monetary policy $\Rightarrow$ limited success; high inflation (10--20\%)
		\item \textbf{Stagflation}: simultaneous stagnation and inflation
	\end{itemize}
\end{frame}

\begin{frame}{Search Theory and the Lucas Critique}
	\textbf{Search models} (McCall, 1970; Mortensen, 1972):
		\vspace{-6pt}
	\begin{itemize}
		\item Developed in early 1970s, before the crisis
		\item Became cornerstone of unemployment theory (Chapter 20)
		\item Generated new data needs (e.g., time spent job searching, JOLTS)
	\end{itemize}

	\textbf{The Lucas Critique}:
		\vspace{-6pt}
	\begin{itemize}
		\item Reduced-form relationships between aggregates can \textbf{break down} if policy changes
		\item Case in point: the \textbf{Phillips curve} ($u$ vs.\ $\pi$)
		\item Lucas (1972): exploiting this trade-off with monetary policy would make the relationship itself break down
		\item When it did break down in the 1970s, Lucas's critique gained added force
		\item Marked a clear break in the development of macroeconomics
	\end{itemize}
\end{frame}

\begin{frame}{Labor Market Data}
	Much of our labor market information comes from \textbf{surveys}:
	\vspace{-6pt}
	\begin{itemize}
		\item \textbf{CPS} (Current Population Survey): monthly, BLS + Census Bureau; measures employment, unemployment, hours, wages; limited panel feature
		\item \textbf{PSID} (Panel Study of Income Dynamics): University of Michigan since 1968; follows individuals over time
		\item \textbf{NLS, SIPP}: other individual panel datasets (BLS, Census Bureau)
		\item \textbf{JOLTS} (Job Openings and Labor Turnover Survey): BLS; data on labor market flows (vacancies, hires, separations); invaluable for search models
		\item \textbf{ATUS} (American Time Use Survey): BLS since 2003; detailed daily activities including job search time
		\item \textbf{Firm-level surveys}: Annual Survey of Manufactures (Census Bureau); report work hours for employees
	\end{itemize}
\end{frame}

\begin{frame}{The 2021--2022 Inflation Echo}
	\textbf{2021--2022}: Sharp rise in inflation, amplified by Russia-Ukraine war energy shock. Fears of a new stagflation period.
	
	\bigskip
	\textbf{Response}: Central banks raised interest rates sharply (2022--2023).
	
	\bigskip
	\textbf{Outcome}: Inflation fell back (energy costs eased, supply bottlenecks decompressed), labor markets stayed tight. Feared stagflation did not materialize.
	
	\bigskip
	\textbf{Lesson}: Macroeconomists are better equipped today, but every challenge has unique properties.
\end{frame}


% ============================================================
\section{Kydland and Prescott: Quantitative Theory}
% ============================================================

\begin{frame}{Kydland and Prescott (1982)}
	Lucas's critique was not just destructive---he suggested building on \textbf{microeconomic theory}.
	
	\bigskip
	\textbf{Kydland and Prescott (1982)} showed how to do this with \textbf{quantitative theory}:
	\begin{itemize}
		\item Base microeconomics on empirical studies in applied fields (labor economics, consumption)
		\item Import parameter values from empirical microeconomic literatures
		\item Compare model output systematically with data
		\item Data must be \textbf{filtered} before comparison (separate trend from cycle)
	\end{itemize}
	
	\bigskip
	\textbf{Result}: An explosion in macroeconomic studies.
\end{frame}
\begin{frame}{Filtering}
	Macroeconomic models define stochastic processes. To compare with data, must extract the aspect of the data under study.
	
	\bigskip
	\begin{itemize}
		\item Stochastic processes can be thought of as sums of sub-processes at different frequencies
		\item Business cycles: periodicities of $\sim$1.5--8 years
		\item \textbf{Band-pass filters}: remove frequencies outside a specified range
		\item Kydland and Prescott used the \textbf{Hodrick-Prescott (HP) filter}---intuitive, widely used
		\item Medium/long-run studies: remove high frequencies, retain low/medium
		\item Financial economics: retain only very high frequencies
	\end{itemize}
\end{frame}
\begin{frame}{From Perfectly Working Markets to New Keynesian Models}
	\textbf{First wave}: perfectly working markets, no monetary policy (``real business cycle'' models).
	
	\textbf{Subsequent enrichments}:
	\vspace{-6pt}
	\begin{itemize}
		\item Added frictions: virtually no practical macro model today is free of market imperfections
		\item Key friction: \textbf{price stickiness} (firms face costs in adjusting prices, adjust infrequently)
		\item Makes monetary policy have direct effects on the economy
		\item $\Rightarrow$ \textbf{New Keynesian framework}: microfounded + sticky prices (and/or sticky wages)
	\end{itemize}
	

	\textbf{Measurement response}: Bils and Klenow (2004) studied CPI survey data, recorded the frequency of price adjustments $\Rightarrow$ parameterized adjustment-cost models.
\end{frame}

% ============================================================
\section{Waves of Macroeconometrics}
% ============================================================

\begin{frame}{Comparing Models to Data: A Methodological Challenge}
	\textbf{``All models are wrong but some can be useful.''}
	
	\bigskip
	A useful model bridges the gap from observable data to the question we want to answer. Knowing which data are important to match is ``often argued to be an art. But art, too, can be taught.''
	
	\bigskip
	\textbf{Evolution of methods}:
	\begin{itemize}
		\item \textbf{Keynesian era}: Large estimated systems with ad-hoc dynamics, lags for better fit
		\item \textbf{Post-Lucas critique}: Back to the drawing board
		\begin{itemize}
			\item Structural estimation: maximum likelihood, Bayesian techniques
			\item Generalized Method of Moments (GMM)
		\end{itemize}
	\end{itemize}
\end{frame}


\begin{frame}{Calibration}
	\textbf{Calibration} (Kydland and Prescott): Select model parameters based on moments of the data \emph{other} than those in focus.
	

	Two sources of parameter values:
	\vspace{-6pt}
	\begin{enumerate}
		\item \textbf{Micro data}: risk aversion, intertemporal substitution
		\item \textbf{Long-run facts}: low-frequency data not in focus for high-frequency questions
	\end{enumerate}
	
	Then derive the model's predictions for the phenomenon of interest.

	\textbf{Spirit}: Parameters are not chosen to match the target moments. Broad patterns matter more than precise fit. Can the model ``at all account for the phenomenon, or are the magnitudes severely off?''
	
	\textbf{Limitation}: Does not lend itself to formal hypothesis testing (shared with Bayesian analysis).
\end{frame}

\begin{frame}{VARs and Natural Experiments}
	\textbf{VARs} (Sims, 1980): A much less structural approach.
		\vspace{-6pt}
	\begin{itemize}
		\item Focus on identifying causal effects in aggregate data
		\item Observe plausibly exogenous shocks 
		\item Trace out impulse response functions
		\item Very common tool in macroeconomics
	\end{itemize}
	
	\textbf{Natural experiments}: Increasingly used to identify causal relationships at the micro level.
		\vspace{-6pt}
	\begin{itemize}
		\item Johnson, Parker, and Souleles (2006): random timing of US tax rebates $\Rightarrow$ measure MPC
		\item Acemoglu, Johnson, and Robinson (2001): colonial-era institutions $\Rightarrow$ long-run growth
		\item Romer and Romer (1989): narrative approach to monetary policy (Fed meeting minutes)
		\item Big-data tools and text analysis entering the field
	\end{itemize}
\end{frame}


\begin{frame}{Intuition vs.\ Quantitative Use}
	Kydland and Prescott also changed how macroeconomists build models:
	
	\bigskip
	\begin{itemize}
		\item \textbf{Before}: Models aimed at building intuitive understanding of mechanisms (e.g., Lucas 1972)
		\item \textbf{After}: Larger, parameterized models delivering quantitative output
		\begin{itemize}
			\item Aggregate across many mechanisms within a single model
			\item Often nonlinear, impossible with pencil and paper
			\item Solve numerically, simulate, compare model output (Chapter 10)
		\end{itemize}
	\end{itemize}
	
	\bigskip
	Has \textbf{not} replaced the need for small models to build intuition. But the aim is to ultimately have both intuitive understanding \emph{and} quantitative assessments.
\end{frame}

% ============================================================
\section{Macroeconomics and Inequality}
% ============================================================

\begin{frame}{Wage and Income Inequality}
	Late 1970s: sharply rising wage inequality, continuing (with varying intensity) to the present.
	

	\textbf{Candidate explanations}:
		\vspace{-6pt}
	\begin{itemize}
		\item Technological change (skill-biased, automation, AI)
		\item Increased exposure to trade
		\item Changes in unionization
	\end{itemize}
	

	\textbf{Labor share}: Remarkably stable during postwar US, but has trended downward in recent decades (also in other countries). Hypotheses: structural change, technological change, changes in competition.
	

	\textbf{Additional dimensions}: Gender wage gaps, racial and ethnic gaps, relative hours worked across groups.
\end{frame}


\begin{frame}{Wealth Inequality and Heterogeneous-Agent Models}
	\textbf{Why macroeconomists care about wealth inequality}:
		\vspace{-6pt}
	\begin{enumerate}
		\item \textbf{Intrinsic interest}: inequality is undesirable; political stability concerns (Piketty, 2014)
		\item \textbf{Heterogeneity matters for policy}: cash-constrained households spend tax rebates; rich households save them $\Rightarrow$ distributional data determines efficacy of interventions
	\end{enumerate}
	
	\textbf{Heterogeneous-agent models} (began 1990s):
		\vspace{-6pt}
	\begin{itemize}
		\item Response to both reasons above
		\item Now a major workhorse for business cycles, inequality, and government policy
		\item Generated a very large literature
	\end{itemize}
	
	
	\textbf{Data}: Various sources indicate significant increase in US wealth inequality since late 1970s. Similar trends documented in other countries.
\end{frame}

\begin{frame}{Microeconomic and Administrative Data}
	\begin{itemize}
		\item \textbf{Survey data}: CPS, PSID, SCF---long the main source of micro data
		\item \textbf{Administrative data} (tax records, government programs):
		\begin{itemize}
			\item More accurate than self-reported surveys
			\item Often much larger (sometimes entire population)
			\item But narrow in scope and tightly controlled for privacy
		\end{itemize}
		\item \textbf{Business data}: payroll companies (e.g., Grigsby, Hurst, Yildirmaz, 2021), custom surveys
		\item \textbf{Matched employer-employee data}: which firms match with which workers; how wages are set
	\end{itemize}
	
	\bigskip
	\textbf{Firm data}: Census Bureau, BLS, commercial vendors (Compustat, Orbis, NETS). Includes entry/exit dynamics, R\&D expenditures, patent data.
\end{frame}

\begin{frame}{Measuring Individual Wealth}
	\begin{itemize}
		\item US does not tax wealth $\Rightarrow$ no direct administrative data
		\item \textbf{Survey of Consumer Finances (SCF)}: Federal Reserve Board, every 3 years since 1983; voluntary but efforts at representativeness
		\item \textbf{Capitalization method}: Use IRS capital income data + assumed rate of return to estimate wealth
		\item \textbf{World Inequality Lab}: Combines national accounts, balance sheets, surveys, tax data, and ``rich lists'' (\url{https://wid.world/})
		\item Small set of countries have taxed wealth and therefore have administrative data
	\end{itemize}
\end{frame}


% ============================================================
\section{Taxes, Government, and Time Consistency}
% ============================================================

\begin{frame}{Taxes and Government Activities}
	Many western countries have experienced slow, long-run increases in the role of government:
			\vspace{-6pt}
	\begin{itemize}
		\item Share of GDP, employment, and transfers (Social Security, welfare)
		\item US marginal tax rates peaked before Reagan, then progressivity lowered; corporate tax rates have also changed
	\end{itemize}
	

	\textbf{Key questions}:
			\vspace{-6pt}
	\begin{itemize}
		\item How have tax changes affected hours worked, economic activity, and inequality?
		\item Choice between different tax bases: capital income, labor income, corporate, sales, property, seignorage
		\item Intersection of macroeconomics and public finance
		\item Public debt: crowding out of private investment, debt sustainability, sovereign stress
	\end{itemize}
\end{frame}


\begin{frame}{Time Inconsistency}
	\textbf{Tax policy}: Efficient to levy taxes on already-installed capital (non-distorting). But if investors know future capital will be taxed, current investment is distorted.
	
	
	$\Rightarrow$ Government would like to promise not to tax future capital, but ex post wants to change its mind: \textbf{time inconsistency problem} (Kydland and Prescott, 1977).
	

	\textbf{Monetary policy}: Similar tensions---promise to be tough on inflation to shape expectations, but temptation to renege.
	

	$\Rightarrow$ \textbf{Central bank independence} and inflation-targeting strategies designed to enhance commitment power. Shaped by lessons of the 1970s--80s.
	
	
	More broadly: interaction between macroeconomics and politics is an active sub-field.
\end{frame}

% ============================================================
\section{The Great Recession}
% ============================================================

\begin{frame}{The Great Moderation and Its End}
	\begin{itemize}
		\item 1979: Volcker's tight monetary policy $\Rightarrow$ recessions in early 1980s, but inflation came down
		\item 1999: Euro created (ten currencies replaced)
		\item Mid-1980s to 2008: The \textbf{Great Moderation}---healthy growth, low volatility
		\begin{itemize}
			\item Attributed to: transparent Fed policy, changes in technological shocks, New Keynesian framework at central banks
		\end{itemize}
	\end{itemize}
	
	\textbf{Then, 2007}: A severe, worldwide economic downturn.
		\vspace{-6pt}
	\begin{itemize}
		\item Not as deep as the Great Depression, but: unemployment rose sharply, recovered very slowly
		\item 2009: Eurozone debt crisis (Greece, Cyprus, Ireland, Portugal, Spain)
		\item Significant uncertainty about government policy paths
	\end{itemize}
\end{frame}

\begin{frame}{Causes and Consequences}
	\textbf{Consensus}: Combination of excessive risk-taking in housing markets + severe frictions in financial markets.
	
	\textbf{Research response}:
		\vspace{-6pt}
	\begin{itemize}
		\item Research on financial markets, institutions, and regulation
		\item Asset markets, debt buildups, and risk-taking intricately intertwined with macroeconomy
		\item Downward trend in real interest rates $\to$ more severe asset-price fluctuations, bubbles, cryptocurrency
		\item Financial frictions now play a prominent role in macro models
	\end{itemize}
	

	\textbf{Basic framework not abandoned}, but enriched in the direction of financial frictions.
	

	\textbf{``The next recession will rarely have the same characteristics as the most recent one''}---theories grow richer, rather than cycling back and forth.
\end{frame}

\begin{frame}{Financial Data}
	\begin{itemize}
		\item \textbf{Key question}: Financial stability---domino effects from interdependencies in portfolios and liabilities
		\item The Great Recession: small mortgage-market shock had outsized effects due to opaque packaging and distribution of mortgage liabilities
		\item \textbf{Flow of Funds} (Federal Reserve Board): quarterly Financial Accounts of the US---assets and liabilities of households, businesses, governments, financial institutions
		\begin{itemize}
			\item Aggregate data; useful for tracking indebtedness trends
			\item Not quite sufficient for detecting inter-institution interdependencies
		\end{itemize}
		\item Asset prices for publicly traded firms: widely available
		\item Non-traded firms: much harder (private equity, venture capital increasingly important)
		\item Individual portfolio decisions: limited to surveys (e.g., SCF)
	\end{itemize}
\end{frame}

% ============================================================
\section{Global Interactions}
% ============================================================

\begin{frame}{International Transmission and Sovereign Default}
	\textbf{The Great Recession was global}:
	\vspace{-6pt}
	\begin{itemize}
		\item Stresses from US housing/banking spread through funding markets, cross-border exposures, trade collapse
		\item Highlighted importance of open-economy models
		\item Shocks propagate via capital flows and relative price movements (terms of trade, real exchange rates)
	\end{itemize}
	

	\textbf{European sovereign-debt crisis (2011--12)}:
	\vspace{-6pt}
	\begin{itemize}
		\item Exposed fragilities in euro-area financial system
		\item Greece, Cyprus restructured debt; Ireland, Portugal, Spain saw spreads spike
		\item Brought sovereign-default models (originally for Latin American crises of the 1980s) back to center stage
	\end{itemize}
	

	\textbf{Post-pandemic}: Rising public debt ratios worldwide keep these topics salient.
\end{frame}

% ============================================================
\section{Climate Change and Energy Economics}
% ============================================================

\begin{frame}{Climate Change and Energy}
	\textbf{Climate change}: Intersection of macroeconomics and environmental economics was near-empty until late 20th century.
	
	\textbf{Macroeconomists' contributions}:
		\vspace{-6pt}
	\begin{itemize}
		\item \textbf{Damage measurement}: studying how aggregates react to weather and climate
		\item \textbf{Integrated assessment models}: examining how different policies influence both economic and climate outcomes jointly
		\item Upsurge in quantitatively oriented modeling of global macro-climate interactions
	\end{itemize}
	

	\textbf{Energy economics}: Attention from the 1970s oil shocks, underscored by recent events (Ukraine war).
		\vspace{-6pt}
	\begin{itemize}
		\item Nature of energy supply; how the economy reacts to shortages
		\item Potential for technological change in energy
		\item Research result communication to policymakers is valuable
	\end{itemize}
\end{frame}

% ============================================================
\section{Where Do We Stand?}
% ============================================================

\begin{frame}{Self-Criticism and Progress}
	\textbf{Question}: Are macroeconomists better positioned today to meet the next challenge?
	
	\begin{itemize}
		\item Macroeconomists {do} admit mistakes---multiple changes in thinking described above
		\item Sometimes theories lacked relevant features $\Rightarrow$ added as swiftly as possible
		\item Sometimes surprised by non-economic shocks that generated economic downturns
	\end{itemize}
	

	Guided by research, macroeconomic policy has been {increasingly successful} over time.
	\vspace{-6pt}
	\begin{itemize}
		\item Responses to the Great Recession and 2020 pandemic benefited in major ways from macroeconomic research
		\item Lessons from one recession are typically useful in the future
	\end{itemize}
\end{frame}

\begin{frame}{Combining Insights, Not Cycling}
	\textbf{We do not cycle back and forth between models.} Rather, we combine insights as they arrive.
	

	\textbf{Example}: Keynes's insights were crucial, but the 1970s brought a break in foundations. Keynesian insights have since been \textbf{added back} into newer, more sophisticated models. These models are clearer on the circumstances under which Keynes's ideas apply.
	

	\textbf{Ongoing developments}:
		\vspace{-6pt}
	\begin{itemize}
		\item Rational expectations may evolve as we get more data on how forecasts are actually made
		\item Financial frictions are now recognized as potentially critical
		\item How to identify key frictions and prevent crises is still open
		\item Rich models are not an end in themselves---simplification for communication and teaching remains essential (IS-LM, RBC models can still build understanding)
	\end{itemize}
\end{frame}
\begin{frame}{Looking Ahead}
	\textbf{Next}: Chapter 2 dives into macroeconomic data and introduces Solow's framework.
	

	The \textbf{neoclassical growth framework} is the core setting used in macroeconomics:
	\vspace{-6pt}
	\begin{itemize}
		\item Precisely motivated by long-run facts that are hard to explain without this theory
		\item Parameters Solow treated as exogenous (saving rates, hours) are better understood as conscious choices in a market environment
	\end{itemize}
	

	Virtually every subsequent chapter builds directly on the market version of the neoclassical growth model:
	\vspace{-6pt}
	\begin{itemize}
		\item Growth, business cycles, asset prices, labor markets, inequality, fiscal and monetary policy, international macro, climate\ldots
	\end{itemize}
	

\end{frame}

\end{document}

